\documentclass[10pt,a4paper]{amsart}
\usepackage[margin=19mm]{geometry}
\usepackage{amsmath,amssymb,mathtools}
\usepackage[scr=rsfs,cal=euler]{mathalfa}
\usepackage{xfrac}
\usepackage[unicode,hidelinks,bookmarks=false]{hyperref}
\newcommand{\ie}{i.e.\ }
\newcommand{\cf}{cf.\ }
\newcommand{\resp}{respectively\ }
\newcommand{\Subsection}{\subsection}
\providecommand{\nobreakdash}{\nobreak\hbox{-}}
\setlength{\emergencystretch}{2em}
\multlinegap0pt
\usepackage{amssymb}
\usepackage{xcolor}
\usepackage{algorithm}\usepackage{algpseudocode}
\usepackage{mathtools}
\newtheorem{prop}{Proposition}
\newcommand{\QQ}{\mathbb{Q}}
\title{Special values of $p$-adic $L$-functions and Iwasawa $\lambda$-invariants of Dirichlet characters}
\author{Heiko Knospe}
\date{Source: arXiv:2401.06100v3 (CC BY 4.0)}
\begin{document}
\maketitle

\noindent\emph{Source and licence.} Special values of $p$-adic $L$-functions and Iwasawa $\lambda$-invariants of Dirichlet characters. Version arXiv:2401.06100v3 (CC BY 4.0), distributed by arXiv under the Creative Commons Attribution 4.0 International licence, http://creativecommons.org/licenses/by/4.0/, as recorded in the arXiv licence metadata for this exact version and shown on the abstract page https://arxiv.org/abs/2401.06100v3. Full licence notice: \texttt{licenses/arxiv-2401.06100-CC-BY-4.0.txt}.\par\smallskip

\noindent\textit{Editorial scope.} Counted unit: the article's prop secondkind (source0.tex lines 567--588). The article's complete original proof of it is delivered with it (source0.tex lines 590--598, one \texttt{proof} environment of 9 lines). No mathematics is written, repaired or re-derived: the statement and the proof are the article's own lines, in the article's own notation, and no retained line is substituted or re-flowed.  No further source-local context is needed: every cross-reference of every retained line resolves inside the two delivered spans.  Nothing was omitted from any retained span, so the package carries no omission record.  No retained line contains a citation, so the wrapper carries no bibliography and no bibliography entry.  No retained line contains a figure environment, an image-inclusion command or drawing code, so no image is delivered and the package declares no asset dependency.  Editorial formatting only, in addition: an AMS \texttt{amsart} preamble that restates the article's own declarations and macros used by the retained lines, namely the environments \texttt{prop} on separate counters, together with the standard packages those lines need.  The retained environments print in this document as Proposition 1 in order of appearance, whereas the article prints the same items with the numbering of its own full document.  The complete native source of the article as distributed in the arXiv e-print for this exact version is bundled unchanged as \texttt{sources/152800/source0.tex}.
\begin{prop} Let $K/\QQ_p$ be a finite extension with ring of integers $\mathcal{O}_K$ and ramification index $e$ at $p$. Let $n$ be a positive integer  and let 
$\pi_n $ be a uniformizer of $\QQ_p(\zeta_{p^n})$, 
e.g.,  $\pi_n = \zeta_{p^n} - 1$. 
Let $t_0 \in \mathcal{O}_K$ with $v_p(t_0) \geq 1$, e.g., $t_0=0$. 
Suppose that
$\mathcal{F}(T)  \in \mathcal{O}_K[[T]]$ is a power series satisfying $\mu_p(\mathcal{F})=0$. 
\begin{enumerate}
\item[(i)] We have $\lambda_p(\mathcal{F}) < \frac{(p-1)p^{n-1}}{e}$ if and only if  $v_p(\mathcal{F}(\pi_n )) < \frac{1}{e}$, and in this case
$$ \lambda_p(\mathcal{F}) = v_p(\mathcal{F}(\pi_n ))\cdot (p-1)p^{n-1} .\\ $$ 
\item[(ii)] Suppose that $\lambda_p(\mathcal{F}) > 0$. Then the equivalence 
$$\lambda_p(\mathcal{F}) < \frac{(p-1)p^{n-1}}{e} + 1 \Longleftrightarrow v_p(\mathcal{F}(\pi_n ) - \mathcal{F}(t_0)) <  \frac{1}{e} + \frac{1}{(p-1)p^{n-1}}$$
holds. In this case,
$$  \lambda_p(\mathcal{F})  = v_p(\mathcal{F}(\pi_n ) - \mathcal{F}(t_0)) \cdot (p-1)p^{n-1}   .$$
\item[(iii)] Suppose that $\lambda_p(\mathcal{F}) > 1$. Then  the equivalence 
$$\lambda_p(\mathcal{F}) < \frac{(p-1)p^{n-1}}{e} + 2 \Longleftrightarrow v_p(\mathcal{F}(\pi_n ) - \mathcal{F}(t_0) -   \mathcal{F}'(t_0)(\pi_n - t_0)) <  \frac{1}{e} + \frac{2}{(p-1)p^{n-1}}$$
holds. In this case,
$$  \lambda_p(\mathcal{F})  = v_p(\mathcal{F}(\pi_n ) - \mathcal{F}(t_0) -   \mathcal{F}'(t_0 ) (\pi_n - t_0) ) \cdot (p-1)p^{n-1}  .$$

\end{enumerate}

\label{secondkind}
\end{prop}

\begin{proof} The statement is true for $\lambda= \lambda_p(\mathcal{F})=0$, and we can assume that $\lambda   \geq 1$. 
Let $\mathcal{F}(T)=a_0 + a_1 (T-t_0) + \dots + a_{\lambda-1}(T-t_0)^{\lambda -1} + a_{\lambda} (T-t_0)^{\lambda} + \dots $. 
We extend the valuation $v_p$ to all finite extensions of $K$ inside a fixed algebraic closure. 
Since $v_p(t_0) > 0$,  
the characterization of $\lambda$ via valuations of the coefficients $a_j$ does not depend on the chosen center $t_0$.
Furthermore, $v_p(a_0), \dots , v_p(a_{\lambda-1}) \geq  \frac{1}{e}$ and $v_p(a_{\lambda})=0$. By assumption, $v_p(\pi_n) = v_p(\pi_n - t_0) = \frac{1}{(p-1)p^{n-1}}$. If $\lambda <  \frac{(p-1)p^{n-1}}{e}$ we have $v_p(\mathcal{F}(\pi_n)) = \frac{\lambda}{(p-1)p^{n-1}} < \frac{1}{e}$. Conversely, suppose that $\lambda \geq \frac{(p-1)p^{n-1}}{e}$.
Then  $v_p(\mathcal{F}(\pi_n )) \geq \frac{1}{e}$. 
This proves (i), and parts (ii) and (iii) can be shown similarly by subtracting the terms of degree $0$ and $1$  from $\mathcal{F}(T)$.
\end{proof}

\end{document}
